\section{Results}
\label{sec:results}

\subsection{Broad Retrieval Gains with a Compact Adapter}
% Generated by paper/scripts/build_evidence.py; do not hand-edit.
\begin{table}[t]
\centering
\small
\setlength{\tabcolsep}{4pt}
\caption{Fixed-anchor retrieval quality, equally weighted over four development collections, two migrations and three budgets. All methods fit the same 1,024 pairs. Differences are ARC2 minus baseline; 99\% individual intervals give approximate 95\% coverage for the five-comparison family.}\label{tab:dev}
\begin{tabular}{lrrrr}
\toprule
Method & nDCG@10 & R@100 & $\Delta$ nDCG & Family interval \\
\midrule
\textbf{ARC2-Global} & 0.5275 & 0.8659 & --- & --- \\
Procrustes & 0.5032 & 0.8661 & +0.0243 & [+0.0181, +0.0308] \\
LA2M20 & 0.4234 & 0.7758 & +0.1041 & [+0.0861, +0.1229] \\
SimpleLinear & 0.3554 & 0.6130 & +0.1721 & [+0.1499, +0.1947] \\
Affine ridge & 0.4318 & 0.7441 & +0.0957 & [+0.0754, +0.1164] \\
Converter & 0.3663 & 0.7641 & +0.1612 & [+0.1427, +0.1799] \\
\bottomrule
\end{tabular}
\end{table}

\begin{figure}[t]
  \centering
  \includegraphics[width=\textwidth]{dev_quality.pdf}
  \caption{Development quality versus anchor-inclusive paid-document coverage.
  Each point equally averages the four collection means after averaging five
  mapping seeds. Native references use the same collection weighting.
  Curves are descriptive; no pooled uncertainty band is implied.}
  \label{fig:dev-quality}
\end{figure}

\method{} delivers the highest macro \ndcg{} among all six evaluated
methods at matched fitting exposure (Table~\ref{tab:dev}). Against the
closest comparator, Procrustes, the gain is $0.024325$, with family interval
$[0.018132,0.030792]$. The gain appears in both migrations: the descriptive
pair-wise macro differences are $0.025733$ for MiniLM and $0.022917$ for E5.
Across individual conditions, ARC2 has higher nDCG point means in 23 of 24
comparisons with Procrustes and all 24 comparisons with each other baseline.
The breadth of the point gains complements the family-adjusted macro
comparison. Figure~\ref{fig:dev-quality} shows the complete budget curves.

The advantage is concentrated in the top-ranked results: macro Recall is
$0.8659$, versus Procrustes' $0.8661$, with a descriptive difference interval
$[-0.004944,0.004336]$. The exception to the 23-of-24 nDCG point gains is
FiQA/E5 at 75\% ($-0.003583$). Native Full-X and Full-Y remain visible as
full-corpus reference operating points.

\subsection{The Second Correction Adds Measurable Value}
\begin{figure}[t]
  \centering
  \includegraphics[width=\textwidth]{dev_ablation.pdf}
  \caption{Development mechanism controls. Dots are equal-condition macro
  nDCG differences; bars are the saved 99\% intervals of the separate
  five-contrast family. The final two rows vary the penalty while retaining
  two steps.}
  \label{fig:ablation}
\end{figure}

The second update provides a measurable benefit: two versus one step improves
macro nDCG by $0.006924$, with family interval $[0.003809,0.009959]$.
The gain over affine ridge is $0.095653$ for the complete initialized
refinement procedure in this supervision regime (Fig.~\ref{fig:ablation});
this contrast evaluates tangent targeting and the retained alignment prior
together. The penalty controls measure sensitivity around the fixed
setting. Changing $\lambda$ from 10 to 1
gives $+0.003197$, with an interval crossing zero; changing it to 100 gives
$-0.009255$, with a wholly negative interval. We retain the frozen setting.

\subsection{Prospective Gains Transfer to TREC-COVID}
% Generated by paper/scripts/build_evidence.py; do not hand-edit.
\begin{table}[t]
\centering
\small
\setlength{\tabcolsep}{4pt}
\caption{Prospective TREC-COVID nDCG@10, averaged over five mapping seeds. A baseline cell marked $\dagger$ has a positive ARC2-minus-baseline interval in the predeclared 30-contrast family (approximate 95\% family coverage). Unmarked cells have intervals crossing zero. Full-X/Y are native full-corpus references.}\label{tab:trec}
\begin{tabular}{lrrrrrr}
\toprule
Method & \multicolumn{3}{c}{MiniLM $\rightarrow$ DistilRoBERTa} & \multicolumn{3}{c}{E5-small $\rightarrow$ E5-base} \\
\midrule
 & 25\% & 50\% & 75\% & 25\% & 50\% & 75\% \\
\textbf{ARC2-Global} & 0.4619 & 0.5000 & 0.5195 & 0.5628 & 0.5904 & 0.6098 \\
Procrustes & 0.4620 & 0.5046 & 0.5195 & 0.5660 & 0.6067 & 0.6198 \\
LA2M20 & 0.4634 & 0.5073 & 0.5186 & 0.5206 & 0.5898 & 0.6188 \\
SimpleLinear & 0.4610 & 0.5076 & 0.5193 & 0.4299$^{\dagger}$ & 0.5287$^{\dagger}$ & 0.5785 \\
Affine ridge & 0.4322 & 0.4854 & 0.5012$^{\dagger}$ & 0.2986$^{\dagger}$ & 0.3281$^{\dagger}$ & 0.3278$^{\dagger}$ \\
Converter & 0.4224$^{\dagger}$ & 0.4780$^{\dagger}$ & 0.5033 & 0.3056$^{\dagger}$ & 0.3590$^{\dagger}$ & 0.4098$^{\dagger}$ \\
Native Full-X & \multicolumn{3}{c}{0.4935 (native)} & \multicolumn{3}{c}{0.6805 (native)} \\
Native Full-Y & \multicolumn{3}{c}{0.5203 (native)} & \multicolumn{3}{c}{0.6421 (native)} \\
\bottomrule
\end{tabular}
\end{table}

\begin{figure}[t]
  \centering
  \includegraphics[width=\textwidth]{trec_quality.pdf}
  \caption{Prospective TREC-COVID quality at fixed 1,024-anchor supervision.
  All six methods and native references are shown; points average the five
  mapping seeds. The two panels use different vertical ranges.}
  \label{fig:trec-quality}
\end{figure}

The held collection confirms a substantial E5 advantage over the evaluated
Converter and affine ridge at \textbf{all three migration budgets}.
At 25\% destination coverage, the gain over Converter is $0.257234$, with
family interval $[0.175731,0.337569]$. At 50\%, ARC2 reaches $0.5904$ nDCG,
versus $0.3590$ for Converter and $0.3281$ for ridge (Table~\ref{tab:trec}).
Across both migrations, family-adjusted improvements occur in five of six
Converter comparisons, four ridge comparisons and two SimpleLinear
comparisons: \textbf{11 improvements in the 30-contrast family}.
Table~\ref{tab:trec} marks every supported pairwise improvement and retains
all six methods at all three budgets.
Figure~\ref{fig:trec-quality} shows how the gains change with coverage.

The MiniLM migration supplies a useful partial-coverage operating point:
at 75\% coverage, \method\ reaches $0.5195$ nDCG against native Full-Y's
$0.5203$. These point estimates illustrate proximity to destination-model
quality while one quarter of the gallery still uses proxies. The
predeclared inferential family compares adapters; it does not establish
equivalence to native Full-Y.

The strongest geometric alternatives remain competitive: Procrustes has
higher point means in all six cells, while the Procrustes and LA2M family
intervals cross zero. Thus the prospective advantage is specific to the
evaluated Converter, ridge and SimpleLinear configurations. The native
references in Table~\ref{tab:trec} retain the quality context for cutover.

\subsection{Using Paid Pairs Strengthens MiniLM Adaptation}
% Generated by paper/scripts/build_evidence.py; do not hand-edit.
\begin{table}[t]
\centering
\small
\setlength{\tabcolsep}{4pt}
\caption{FiQA at the same 50\% paid-document budget. Fixed uses 1,024 fitting pairs; All uses all 28,819 paid pairs. N denotes nDCG@10 and R denotes Recall@100. Native references use their own query/document encoders.}\label{tab:allpaid}
\begin{tabular}{lrrrrrr}
\toprule
Method & \multicolumn{3}{c}{MiniLM $\rightarrow$ DistilRoBERTa} & \multicolumn{3}{c}{E5-small $\rightarrow$ E5-base} \\
\midrule
 & Fixed N & All N & All R & Fixed N & All N & All R \\
\textbf{ARC2-Global} & 0.3075 & 0.3671 & 0.6922 & 0.3116 & 0.2919 & 0.5595 \\
Procrustes & 0.2858 & 0.3334 & 0.6905 & 0.3100 & 0.3362 & 0.6534 \\
LA2M20 & 0.2377 & 0.3389 & 0.6950 & 0.2300 & 0.3293 & 0.6473 \\
SimpleLinear & 0.2360 & 0.3323 & 0.6778 & 0.1942 & 0.1679 & 0.3618 \\
Affine ridge & 0.2890 & 0.3673 & 0.6800 & 0.1479 & 0.2270 & 0.4245 \\
Converter & 0.2496 & 0.3258 & 0.6633 & 0.1397 & 0.2625 & 0.5789 \\
Native Full-X & \multicolumn{2}{c}{0.3900 (native)} & 0.7002 & \multicolumn{2}{c}{0.3465 (native)} & 0.6453 \\
Native Full-Y & \multicolumn{2}{c}{0.4047 (native)} & 0.7039 & \multicolumn{2}{c}{0.3891 (native)} & 0.6968 \\
\bottomrule
\end{tabular}
\end{table}

More fitting supervision produces a strong MiniLM operating point without
increasing the paid-document set. With all paid FiQA pairs, ARC2 reaches
$0.367146$ nDCG and exceeds
Procrustes, LA2M, SimpleLinear, and Converter in the locked family comparisons;
ridge remains unresolved. E5 shows a different ordering: ARC2 trails
Procrustes and LA2M, exceeds SimpleLinear and ridge, and is unresolved
against Converter. Table~\ref{tab:allpaid} retains both migrations and all
methods. Fixed-anchor MiniLM ARC2 trails each all-paid baseline in the
locked cross-regime nDCG comparisons. Increasing ARC2's fitting exposure raises MiniLM nDCG from
$0.307505$ to $0.367146$; E5 changes from $0.311635$ to $0.291867$, with
that nDCG interval unresolved. All-paid fitting retains $\lambda=10$ in
the summed objective, so it evaluates the fixed procedure at larger
supervision rather than isolating data quantity at a fixed effective penalty.
The result identifies when the paid pairs
can improve adaptation, while the two-migration comparison captures the change
in method ordering.

\paragraph{A practical insight: validate the mixed gallery.}
Holding each fitted mapper fixed, we compare proxy-only and mixed
galleries under the same document replacement sets.
For all-paid E5 ARC2 on FiQA, nDCG decreases from $0.314908$ to $0.291867$;
Recall decreases from $0.654164$ to $0.559490$. The paid-document share of
top-10 results falls from $0.491640$ to $0.064720$ despite 50\% paid coverage.
For MiniLM ARC2, nDCG instead rises from $0.350951$ to $0.367146$.
These opposite replacement effects make mixed-gallery relevance the
decision metric for partial migration. Good proxy-only quality is
insufficient to select a cutover point; the documents must be ranked
together in the intended gallery.
